% Options for packages loaded elsewhere
% Options for packages loaded elsewhere
\PassOptionsToPackage{unicode}{hyperref}
\PassOptionsToPackage{hyphens}{url}
\PassOptionsToPackage{dvipsnames,svgnames,x11names}{xcolor}
%
\documentclass[
  letterpaper,
  DIV=11,
  numbers=noendperiod]{scrartcl}
\usepackage{xcolor}
\usepackage{amsmath,amssymb}
\setcounter{secnumdepth}{-\maxdimen} % remove section numbering
\usepackage{iftex}
\ifPDFTeX
  \usepackage[T1]{fontenc}
  \usepackage[utf8]{inputenc}
  \usepackage{textcomp} % provide euro and other symbols
\else % if luatex or xetex
  \usepackage{unicode-math} % this also loads fontspec
  \defaultfontfeatures{Scale=MatchLowercase}
  \defaultfontfeatures[\rmfamily]{Ligatures=TeX,Scale=1}
\fi
\usepackage{lmodern}
\ifPDFTeX\else
  % xetex/luatex font selection
\fi
% Use upquote if available, for straight quotes in verbatim environments
\IfFileExists{upquote.sty}{\usepackage{upquote}}{}
\IfFileExists{microtype.sty}{% use microtype if available
  \usepackage[]{microtype}
  \UseMicrotypeSet[protrusion]{basicmath} % disable protrusion for tt fonts
}{}
\makeatletter
\@ifundefined{KOMAClassName}{% if non-KOMA class
  \IfFileExists{parskip.sty}{%
    \usepackage{parskip}
  }{% else
    \setlength{\parindent}{0pt}
    \setlength{\parskip}{6pt plus 2pt minus 1pt}}
}{% if KOMA class
  \KOMAoptions{parskip=half}}
\makeatother
% Make \paragraph and \subparagraph free-standing
\makeatletter
\ifx\paragraph\undefined\else
  \let\oldparagraph\paragraph
  \renewcommand{\paragraph}{
    \@ifstar
      \xxxParagraphStar
      \xxxParagraphNoStar
  }
  \newcommand{\xxxParagraphStar}[1]{\oldparagraph*{#1}\mbox{}}
  \newcommand{\xxxParagraphNoStar}[1]{\oldparagraph{#1}\mbox{}}
\fi
\ifx\subparagraph\undefined\else
  \let\oldsubparagraph\subparagraph
  \renewcommand{\subparagraph}{
    \@ifstar
      \xxxSubParagraphStar
      \xxxSubParagraphNoStar
  }
  \newcommand{\xxxSubParagraphStar}[1]{\oldsubparagraph*{#1}\mbox{}}
  \newcommand{\xxxSubParagraphNoStar}[1]{\oldsubparagraph{#1}\mbox{}}
\fi
\makeatother

\usepackage{color}
\usepackage{fancyvrb}
\newcommand{\VerbBar}{|}
\newcommand{\VERB}{\Verb[commandchars=\\\{\}]}
\DefineVerbatimEnvironment{Highlighting}{Verbatim}{commandchars=\\\{\}}
% Add ',fontsize=\small' for more characters per line
\usepackage{framed}
\definecolor{shadecolor}{RGB}{241,243,245}
\newenvironment{Shaded}{\begin{snugshade}}{\end{snugshade}}
\newcommand{\AlertTok}[1]{\textcolor[rgb]{0.68,0.00,0.00}{#1}}
\newcommand{\AnnotationTok}[1]{\textcolor[rgb]{0.37,0.37,0.37}{#1}}
\newcommand{\AttributeTok}[1]{\textcolor[rgb]{0.40,0.45,0.13}{#1}}
\newcommand{\BaseNTok}[1]{\textcolor[rgb]{0.68,0.00,0.00}{#1}}
\newcommand{\BuiltInTok}[1]{\textcolor[rgb]{0.00,0.23,0.31}{#1}}
\newcommand{\CharTok}[1]{\textcolor[rgb]{0.13,0.47,0.30}{#1}}
\newcommand{\CommentTok}[1]{\textcolor[rgb]{0.37,0.37,0.37}{#1}}
\newcommand{\CommentVarTok}[1]{\textcolor[rgb]{0.37,0.37,0.37}{\textit{#1}}}
\newcommand{\ConstantTok}[1]{\textcolor[rgb]{0.56,0.35,0.01}{#1}}
\newcommand{\ControlFlowTok}[1]{\textcolor[rgb]{0.00,0.23,0.31}{\textbf{#1}}}
\newcommand{\DataTypeTok}[1]{\textcolor[rgb]{0.68,0.00,0.00}{#1}}
\newcommand{\DecValTok}[1]{\textcolor[rgb]{0.68,0.00,0.00}{#1}}
\newcommand{\DocumentationTok}[1]{\textcolor[rgb]{0.37,0.37,0.37}{\textit{#1}}}
\newcommand{\ErrorTok}[1]{\textcolor[rgb]{0.68,0.00,0.00}{#1}}
\newcommand{\ExtensionTok}[1]{\textcolor[rgb]{0.00,0.23,0.31}{#1}}
\newcommand{\FloatTok}[1]{\textcolor[rgb]{0.68,0.00,0.00}{#1}}
\newcommand{\FunctionTok}[1]{\textcolor[rgb]{0.28,0.35,0.67}{#1}}
\newcommand{\ImportTok}[1]{\textcolor[rgb]{0.00,0.46,0.62}{#1}}
\newcommand{\InformationTok}[1]{\textcolor[rgb]{0.37,0.37,0.37}{#1}}
\newcommand{\KeywordTok}[1]{\textcolor[rgb]{0.00,0.23,0.31}{\textbf{#1}}}
\newcommand{\NormalTok}[1]{\textcolor[rgb]{0.00,0.23,0.31}{#1}}
\newcommand{\OperatorTok}[1]{\textcolor[rgb]{0.37,0.37,0.37}{#1}}
\newcommand{\OtherTok}[1]{\textcolor[rgb]{0.00,0.23,0.31}{#1}}
\newcommand{\PreprocessorTok}[1]{\textcolor[rgb]{0.68,0.00,0.00}{#1}}
\newcommand{\RegionMarkerTok}[1]{\textcolor[rgb]{0.00,0.23,0.31}{#1}}
\newcommand{\SpecialCharTok}[1]{\textcolor[rgb]{0.37,0.37,0.37}{#1}}
\newcommand{\SpecialStringTok}[1]{\textcolor[rgb]{0.13,0.47,0.30}{#1}}
\newcommand{\StringTok}[1]{\textcolor[rgb]{0.13,0.47,0.30}{#1}}
\newcommand{\VariableTok}[1]{\textcolor[rgb]{0.07,0.07,0.07}{#1}}
\newcommand{\VerbatimStringTok}[1]{\textcolor[rgb]{0.13,0.47,0.30}{#1}}
\newcommand{\WarningTok}[1]{\textcolor[rgb]{0.37,0.37,0.37}{\textit{#1}}}

\usepackage{longtable,booktabs,array}
\newcounter{none} % for unnumbered tables
\usepackage{calc} % for calculating minipage widths
% Correct order of tables after \paragraph or \subparagraph
\usepackage{etoolbox}
\makeatletter
\patchcmd\longtable{\par}{\if@noskipsec\mbox{}\fi\par}{}{}
\makeatother
% Allow footnotes in longtable head/foot
\IfFileExists{footnotehyper.sty}{\usepackage{footnotehyper}}{\usepackage{footnote}}
\makesavenoteenv{longtable}
\usepackage{graphicx}
\makeatletter
\newsavebox\pandoc@box
\newcommand*\pandocbounded[1]{% scales image to fit in text height/width
  \sbox\pandoc@box{#1}%
  \Gscale@div\@tempa{\textheight}{\dimexpr\ht\pandoc@box+\dp\pandoc@box\relax}%
  \Gscale@div\@tempb{\linewidth}{\wd\pandoc@box}%
  \ifdim\@tempb\p@<\@tempa\p@\let\@tempa\@tempb\fi% select the smaller of both
  \ifdim\@tempa\p@<\p@\scalebox{\@tempa}{\usebox\pandoc@box}%
  \else\usebox{\pandoc@box}%
  \fi%
}
% Set default figure placement to htbp
\def\fps@figure{htbp}
\makeatother





\setlength{\emergencystretch}{3em} % prevent overfull lines

\providecommand{\tightlist}{%
  \setlength{\itemsep}{0pt}\setlength{\parskip}{0pt}}



 


\KOMAoption{captions}{tableheading}
\makeatletter
\@ifpackageloaded{caption}{}{\usepackage{caption}}
\AtBeginDocument{%
\ifdefined\contentsname
  \renewcommand*\contentsname{Table of contents}
\else
  \newcommand\contentsname{Table of contents}
\fi
\ifdefined\listfigurename
  \renewcommand*\listfigurename{List of Figures}
\else
  \newcommand\listfigurename{List of Figures}
\fi
\ifdefined\listtablename
  \renewcommand*\listtablename{List of Tables}
\else
  \newcommand\listtablename{List of Tables}
\fi
\ifdefined\figurename
  \renewcommand*\figurename{Figure}
\else
  \newcommand\figurename{Figure}
\fi
\ifdefined\tablename
  \renewcommand*\tablename{Table}
\else
  \newcommand\tablename{Table}
\fi
}
\@ifpackageloaded{float}{}{\usepackage{float}}
\floatstyle{ruled}
\@ifundefined{c@chapter}{\newfloat{codelisting}{h}{lop}}{\newfloat{codelisting}{h}{lop}[chapter]}
\floatname{codelisting}{Listing}
\newcommand*\listoflistings{\listof{codelisting}{List of Listings}}
\makeatother
\makeatletter
\makeatother
\makeatletter
\@ifpackageloaded{caption}{}{\usepackage{caption}}
\@ifpackageloaded{subcaption}{}{\usepackage{subcaption}}
\makeatother
\usepackage{bookmark}
\IfFileExists{xurl.sty}{\usepackage{xurl}}{} % add URL line breaks if available
\urlstyle{same}
% fallback for those not using the hyperref driver hyperxmp:
\makeatletter
\@ifundefined{xmpquote}{\newcommand{\xmpquote}[1]{#1}}{}
\makeatother
\hypersetup{
  pdftitle={Study A (probiotic trial) --- quality control before analysis},
  pdfauthor={Group project},
  colorlinks=true,
  linkcolor={blue},
  filecolor={Maroon},
  citecolor={Blue},
  urlcolor={Blue},
  pdfcreator={LaTeX via pandoc}}


\title{Study A (probiotic trial) --- quality control before analysis}
\author{Group project}
\date{}
\begin{document}
\maketitle


This document checks the \texttt{group\_a\_probiotic} dataset for
problems \textbf{before} we run any tests. QC is about trust: we want to
know the tables line up, the counts look like counts, and nothing
surprising hides in the metadata. We fix nothing here --- we only look
and report.

\begin{Shaded}
\begin{Highlighting}[]
\FunctionTok{library}\NormalTok{(readr)     }\CommentTok{\# reading csv files}
\FunctionTok{library}\NormalTok{(dplyr)     }\CommentTok{\# data wrangling}
\FunctionTok{library}\NormalTok{(tidyr)     }\CommentTok{\# reshaping wide \textless{}{-}\textgreater{} long}
\FunctionTok{library}\NormalTok{(ggplot2)   }\CommentTok{\# plots}

\FunctionTok{source}\NormalTok{(here}\SpecialCharTok{::}\FunctionTok{here}\NormalTok{(}\StringTok{"R/count\_tables.R"}\NormalTok{))}
\NormalTok{d }\OtherTok{\textless{}{-}} \FunctionTok{read\_counts}\NormalTok{(}\StringTok{"group\_a"}\NormalTok{)   }\CommentTok{\# metadata, species, taxonomy, genes, gene\_info, cytokines}
\end{Highlighting}
\end{Shaded}

\subsection{1. Who is in the data?}\label{who-is-in-the-data}

We expect 176 people, sampled pre and post, in two arms.

\begin{Shaded}
\begin{Highlighting}[]
\NormalTok{m }\OtherTok{\textless{}{-}}\NormalTok{ d}\SpecialCharTok{$}\NormalTok{metadata}
\NormalTok{m }\SpecialCharTok{\%\textgreater{}\%} \FunctionTok{count}\NormalTok{(arm, time\_point)          }\CommentTok{\# samples per arm x time point}
\end{Highlighting}
\end{Shaded}

\begin{verbatim}
# A tibble: 4 x 3
  arm       time_point     n
  <chr>     <chr>      <int>
1 placebo   post          55
2 placebo   pre           55
3 probiotic post         121
4 probiotic pre          121
\end{verbatim}

\begin{Shaded}
\begin{Highlighting}[]
\NormalTok{m }\SpecialCharTok{\%\textgreater{}\%} \FunctionTok{distinct}\NormalTok{(participant\_id, arm) }\SpecialCharTok{\%\textgreater{}\%} \FunctionTok{count}\NormalTok{(arm)   }\CommentTok{\# people per arm}
\end{Highlighting}
\end{Shaded}

\begin{verbatim}
# A tibble: 2 x 2
  arm           n
  <chr>     <int>
1 placebo      55
2 probiotic   121
\end{verbatim}

\begin{Shaded}
\begin{Highlighting}[]
\FunctionTok{nrow}\NormalTok{(m); }\FunctionTok{length}\NormalTok{(}\FunctionTok{unique}\NormalTok{(m}\SpecialCharTok{$}\NormalTok{participant\_id))}
\end{Highlighting}
\end{Shaded}

\begin{verbatim}
[1] 352
\end{verbatim}

\begin{verbatim}
[1] 176
\end{verbatim}

\begin{Shaded}
\begin{Highlighting}[]
\CommentTok{\# Does every person have exactly one pre and one post sample?}
\NormalTok{m }\SpecialCharTok{\%\textgreater{}\%} \FunctionTok{count}\NormalTok{(participant\_id) }\SpecialCharTok{\%\textgreater{}\%} \FunctionTok{count}\NormalTok{(n, }\AttributeTok{name =} \StringTok{"people"}\NormalTok{)}
\end{Highlighting}
\end{Shaded}

\begin{verbatim}
# A tibble: 1 x 2
      n people
  <int>  <int>
1     2    176
\end{verbatim}

\subsection{2. Do the tables line up?}\label{do-the-tables-line-up}

Everything joins on \texttt{sample\_id}. \texttt{read\_counts()} already
reorders the count columns to match the metadata rows, but we check it
rather than trust it --- DESeq2 will silently give wrong answers if rows
and columns are mismatched.

\begin{Shaded}
\begin{Highlighting}[]
\FunctionTok{identical}\NormalTok{(}\FunctionTok{colnames}\NormalTok{(d}\SpecialCharTok{$}\NormalTok{species), m}\SpecialCharTok{$}\NormalTok{sample\_id)   }\CommentTok{\# TRUE = order matches}
\end{Highlighting}
\end{Shaded}

\begin{verbatim}
[1] TRUE
\end{verbatim}

\begin{Shaded}
\begin{Highlighting}[]
\FunctionTok{identical}\NormalTok{(}\FunctionTok{colnames}\NormalTok{(d}\SpecialCharTok{$}\NormalTok{genes),   m}\SpecialCharTok{$}\NormalTok{sample\_id)}
\end{Highlighting}
\end{Shaded}

\begin{verbatim}
[1] TRUE
\end{verbatim}

\begin{Shaded}
\begin{Highlighting}[]
\FunctionTok{setequal}\NormalTok{(d}\SpecialCharTok{$}\NormalTok{cytokines}\SpecialCharTok{$}\NormalTok{sample\_id, m}\SpecialCharTok{$}\NormalTok{sample\_id)  }\CommentTok{\# same samples in the cytokine file?}
\end{Highlighting}
\end{Shaded}

\begin{verbatim}
[1] TRUE
\end{verbatim}

\begin{Shaded}
\begin{Highlighting}[]
\CommentTok{\# sample\_id looks like A001\_pre. Check every one follows person + time point.}
\NormalTok{bad\_ids }\OtherTok{\textless{}{-}}\NormalTok{ m }\SpecialCharTok{\%\textgreater{}\%} \FunctionTok{filter}\NormalTok{(}\SpecialCharTok{!}\FunctionTok{grepl}\NormalTok{(}\StringTok{"\^{}(A[0{-}9]\{3\})\_(pre|post)$"}\NormalTok{, sample\_id))}
\FunctionTok{nrow}\NormalTok{(bad\_ids)   }\CommentTok{\# 0 = all ids parse as expected}
\end{Highlighting}
\end{Shaded}

\begin{verbatim}
[1] 0
\end{verbatim}

\subsection{3. Are the count tables really
counts?}\label{are-the-count-tables-really-counts}

Raw counts for DESeq2 must be whole numbers, never negative.

\begin{Shaded}
\begin{Highlighting}[]
\NormalTok{check\_counts }\OtherTok{\textless{}{-}} \ControlFlowTok{function}\NormalTok{(x, label) \{}
  \FunctionTok{cat}\NormalTok{(label, }\StringTok{":}\SpecialCharTok{\textbackslash{}n}\StringTok{"}\NormalTok{,}
      \StringTok{"  smallest value:   "}\NormalTok{, }\FunctionTok{min}\NormalTok{(x), }\StringTok{"}\SpecialCharTok{\textbackslash{}n}\StringTok{"}\NormalTok{,}
      \StringTok{"  any negative:     "}\NormalTok{, }\FunctionTok{any}\NormalTok{(x }\SpecialCharTok{\textless{}} \DecValTok{0}\NormalTok{), }\StringTok{"}\SpecialCharTok{\textbackslash{}n}\StringTok{"}\NormalTok{,}
      \StringTok{"  any non{-}integer:  "}\NormalTok{, }\FunctionTok{any}\NormalTok{(x }\SpecialCharTok{!=} \FunctionTok{round}\NormalTok{(x)), }\StringTok{"}\SpecialCharTok{\textbackslash{}n}\StringTok{"}\NormalTok{,}
      \StringTok{"  any NA:           "}\NormalTok{, }\FunctionTok{anyNA}\NormalTok{(x), }\StringTok{"}\SpecialCharTok{\textbackslash{}n}\StringTok{"}\NormalTok{, }\AttributeTok{sep =} \StringTok{""}\NormalTok{)}
\NormalTok{\}}
\FunctionTok{check\_counts}\NormalTok{(d}\SpecialCharTok{$}\NormalTok{species, }\StringTok{"microbiome (220 species)"}\NormalTok{)}
\end{Highlighting}
\end{Shaded}

\begin{verbatim}
microbiome (220 species):
  smallest value:   0
  any negative:     FALSE
  any non-integer:  FALSE
  any NA:           FALSE
\end{verbatim}

\begin{Shaded}
\begin{Highlighting}[]
\FunctionTok{check\_counts}\NormalTok{(d}\SpecialCharTok{$}\NormalTok{genes,   }\StringTok{"rna{-}seq (14226 genes)"}\NormalTok{)}
\end{Highlighting}
\end{Shaded}

\begin{verbatim}
rna-seq (14226 genes):
  smallest value:   0
  any negative:     FALSE
  any non-integer:  FALSE
  any NA:           FALSE
\end{verbatim}

\subsection{4. Sequencing depth (library
sizes)}\label{sequencing-depth-library-sizes}

The library size is the total count in a sample. Big differences between
samples are normal and DESeq2 handles them --- but samples with absurdly
small totals, or depth that tracks the batch or arm, are worth knowing
about.

\begin{Shaded}
\begin{Highlighting}[]
\NormalTok{m}\SpecialCharTok{$}\NormalTok{lib\_micro }\OtherTok{\textless{}{-}} \FunctionTok{colSums}\NormalTok{(d}\SpecialCharTok{$}\NormalTok{species)}
\NormalTok{m}\SpecialCharTok{$}\NormalTok{lib\_gene  }\OtherTok{\textless{}{-}} \FunctionTok{colSums}\NormalTok{(d}\SpecialCharTok{$}\NormalTok{genes)}
\FunctionTok{summary}\NormalTok{(m[, }\FunctionTok{c}\NormalTok{(}\StringTok{"lib\_micro"}\NormalTok{, }\StringTok{"lib\_gene"}\NormalTok{)])}
\end{Highlighting}
\end{Shaded}

\begin{verbatim}
   lib_micro         lib_gene       
 Min.   :   604   Min.   : 6288748  
 1st Qu.: 23770   1st Qu.:10614512  
 Median : 32749   Median :12624366  
 Mean   : 35435   Mean   :13064536  
 3rd Qu.: 41380   3rd Qu.:14695888  
 Max.   :152613   Max.   :31894266  
\end{verbatim}

\begin{Shaded}
\begin{Highlighting}[]
\NormalTok{arm\_cols }\OtherTok{\textless{}{-}} \FunctionTok{c}\NormalTok{(}\AttributeTok{probiotic =} \StringTok{"\#1B7837"}\NormalTok{, }\AttributeTok{placebo =} \StringTok{"\#E08214"}\NormalTok{)}

\FunctionTok{ggplot}\NormalTok{(m, }\FunctionTok{aes}\NormalTok{(}\AttributeTok{x =}\NormalTok{ lib\_micro, }\AttributeTok{fill =}\NormalTok{ arm)) }\SpecialCharTok{+}
  \FunctionTok{geom\_histogram}\NormalTok{(}\AttributeTok{bins =} \DecValTok{30}\NormalTok{, }\AttributeTok{colour =} \StringTok{"white"}\NormalTok{, }\AttributeTok{linewidth =}\NormalTok{ .}\DecValTok{2}\NormalTok{, }\AttributeTok{alpha =}\NormalTok{ .}\DecValTok{85}\NormalTok{) }\SpecialCharTok{+}
  \FunctionTok{facet\_grid}\NormalTok{(arm }\SpecialCharTok{\textasciitilde{}}\NormalTok{ time\_point) }\SpecialCharTok{+}
  \FunctionTok{scale\_fill\_manual}\NormalTok{(}\AttributeTok{values =}\NormalTok{ arm\_cols) }\SpecialCharTok{+}
  \FunctionTok{labs}\NormalTok{(}\AttributeTok{title =} \StringTok{"Microbiome library sizes"}\NormalTok{, }\AttributeTok{x =} \StringTok{"reads"}\NormalTok{, }\AttributeTok{y =} \StringTok{"samples"}\NormalTok{) }\SpecialCharTok{+}
  \FunctionTok{theme\_minimal}\NormalTok{(}\AttributeTok{base\_size =} \DecValTok{11}\NormalTok{) }\SpecialCharTok{+}
  \FunctionTok{theme}\NormalTok{(}\AttributeTok{legend.position =} \StringTok{"none"}\NormalTok{)}
\end{Highlighting}
\end{Shaded}

\pandocbounded{\includegraphics[keepaspectratio]{01-qc_files/figure-pdf/libsize-plots-1.pdf}}

\begin{Shaded}
\begin{Highlighting}[]
\FunctionTok{ggplot}\NormalTok{(m, }\FunctionTok{aes}\NormalTok{(}\AttributeTok{x =}\NormalTok{ lib\_gene, }\AttributeTok{fill =}\NormalTok{ rnaseq\_batch)) }\SpecialCharTok{+}
  \FunctionTok{geom\_histogram}\NormalTok{(}\AttributeTok{bins =} \DecValTok{30}\NormalTok{, }\AttributeTok{colour =} \StringTok{"white"}\NormalTok{, }\AttributeTok{linewidth =}\NormalTok{ .}\DecValTok{2}\NormalTok{, }\AttributeTok{alpha =}\NormalTok{ .}\DecValTok{85}\NormalTok{) }\SpecialCharTok{+}
  \FunctionTok{facet\_grid}\NormalTok{(arm }\SpecialCharTok{\textasciitilde{}}\NormalTok{ time\_point) }\SpecialCharTok{+}
  \FunctionTok{scale\_fill\_manual}\NormalTok{(}\AttributeTok{values =} \FunctionTok{c}\NormalTok{(}\AttributeTok{B1 =} \StringTok{"\#4393C3"}\NormalTok{, }\AttributeTok{B2 =} \StringTok{"\#D6604D"}\NormalTok{)) }\SpecialCharTok{+}
  \FunctionTok{labs}\NormalTok{(}\AttributeTok{title =} \StringTok{"RNA{-}seq library sizes"}\NormalTok{, }\AttributeTok{fill =} \StringTok{"batch"}\NormalTok{,}
       \AttributeTok{x =} \StringTok{"reads"}\NormalTok{, }\AttributeTok{y =} \StringTok{"samples"}\NormalTok{) }\SpecialCharTok{+}
  \FunctionTok{theme\_minimal}\NormalTok{(}\AttributeTok{base\_size =} \DecValTok{11}\NormalTok{)}
\end{Highlighting}
\end{Shaded}

\pandocbounded{\includegraphics[keepaspectratio]{01-qc_files/figure-pdf/libsize-plots-2.pdf}}

\begin{Shaded}
\begin{Highlighting}[]
\CommentTok{\# Is sequencing batch balanced across arms? If one arm sits mostly in one}
\CommentTok{\# batch, batch effects could masquerade as treatment effects later.}
\FunctionTok{table}\NormalTok{(m}\SpecialCharTok{$}\NormalTok{arm, m}\SpecialCharTok{$}\NormalTok{rnaseq\_batch)}
\end{Highlighting}
\end{Shaded}

\begin{verbatim}
           
             B1  B2
  placebo    70  40
  probiotic 102 140
\end{verbatim}

\subsection{5. Which features are actually
there?}\label{which-features-are-actually-there}

Features that are zero nearly everywhere carry no information; we will
filter them at the DESeq2 step, but here we just count them.

\begin{Shaded}
\begin{Highlighting}[]
\NormalTok{prev\_species }\OtherTok{\textless{}{-}} \FunctionTok{rowSums}\NormalTok{(d}\SpecialCharTok{$}\NormalTok{species }\SpecialCharTok{\textgreater{}} \DecValTok{0}\NormalTok{)}
\NormalTok{prev\_genes   }\OtherTok{\textless{}{-}} \FunctionTok{rowSums}\NormalTok{(d}\SpecialCharTok{$}\NormalTok{genes }\SpecialCharTok{\textgreater{}} \DecValTok{0}\NormalTok{)}
\FunctionTok{cat}\NormalTok{(}\StringTok{"species never seen:          "}\NormalTok{, }\FunctionTok{sum}\NormalTok{(prev\_species }\SpecialCharTok{==} \DecValTok{0}\NormalTok{), }\StringTok{"of"}\NormalTok{, }\FunctionTok{nrow}\NormalTok{(d}\SpecialCharTok{$}\NormalTok{species), }\StringTok{"}\SpecialCharTok{\textbackslash{}n}\StringTok{"}\NormalTok{)}
\end{Highlighting}
\end{Shaded}

\begin{verbatim}
species never seen:           53 of 220 
\end{verbatim}

\begin{Shaded}
\begin{Highlighting}[]
\FunctionTok{cat}\NormalTok{(}\StringTok{"species in \textless{} 10\% of samples: "}\NormalTok{, }\FunctionTok{sum}\NormalTok{(prev\_species }\SpecialCharTok{\textless{}} \FloatTok{0.1} \SpecialCharTok{*} \FunctionTok{ncol}\NormalTok{(d}\SpecialCharTok{$}\NormalTok{species)), }\StringTok{"}\SpecialCharTok{\textbackslash{}n}\StringTok{"}\NormalTok{)}
\end{Highlighting}
\end{Shaded}

\begin{verbatim}
species in < 10% of samples:  158 
\end{verbatim}

\begin{Shaded}
\begin{Highlighting}[]
\FunctionTok{cat}\NormalTok{(}\StringTok{"genes never seen:            "}\NormalTok{, }\FunctionTok{sum}\NormalTok{(prev\_genes }\SpecialCharTok{==} \DecValTok{0}\NormalTok{), }\StringTok{"of"}\NormalTok{, }\FunctionTok{nrow}\NormalTok{(d}\SpecialCharTok{$}\NormalTok{genes), }\StringTok{"}\SpecialCharTok{\textbackslash{}n}\StringTok{"}\NormalTok{)}
\end{Highlighting}
\end{Shaded}

\begin{verbatim}
genes never seen:             0 of 14226 
\end{verbatim}

\begin{Shaded}
\begin{Highlighting}[]
\FunctionTok{cat}\NormalTok{(}\StringTok{"genes in \textless{} 10\% of samples:   "}\NormalTok{, }\FunctionTok{sum}\NormalTok{(prev\_genes }\SpecialCharTok{\textless{}} \FloatTok{0.1} \SpecialCharTok{*} \FunctionTok{ncol}\NormalTok{(d}\SpecialCharTok{$}\NormalTok{genes)), }\StringTok{"}\SpecialCharTok{\textbackslash{}n}\StringTok{"}\NormalTok{)}
\end{Highlighting}
\end{Shaded}

\begin{verbatim}
genes in < 10% of samples:    0 
\end{verbatim}

\begin{Shaded}
\begin{Highlighting}[]
\CommentTok{\# Sanity check against biology: Lactobacillus species should dominate this site.}
\NormalTok{rel }\OtherTok{\textless{}{-}} \FunctionTok{colSums}\NormalTok{(d}\SpecialCharTok{$}\NormalTok{species) }\SpecialCharTok{|\textgreater{}} \FunctionTok{setNames}\NormalTok{(}\ConstantTok{NULL}\NormalTok{)}
\NormalTok{sp\_mean\_rel }\OtherTok{\textless{}{-}} \FunctionTok{sort}\NormalTok{(}\FunctionTok{rowMeans}\NormalTok{(}\FunctionTok{sweep}\NormalTok{(d}\SpecialCharTok{$}\NormalTok{species, }\DecValTok{2}\NormalTok{, rel, }\StringTok{"/"}\NormalTok{)), }\AttributeTok{decreasing =} \ConstantTok{TRUE}\NormalTok{)}
\FunctionTok{head}\NormalTok{(sp\_mean\_rel, }\DecValTok{5}\NormalTok{)}
\end{Highlighting}
\end{Shaded}

\begin{verbatim}
                 lactobacillus_iners    gardnerella_swidsinskii_leopoldii 
                          0.18135638                           0.13985175 
             lactobacillus_crispatus candidatus_lachnocurva_vaginae_bvab1 
                          0.13602666                           0.05582113 
                 fannyhessea_vaginae 
                          0.05566722 
\end{verbatim}

\subsection{6. Metadata gaps and
oddities}\label{metadata-gaps-and-oddities}

\begin{Shaded}
\begin{Highlighting}[]
\CommentTok{\# Missing values per column. "nugent" and friends are clinical measures taken}
\CommentTok{\# at each visit, so some gaps are plausible; batch gaps are not.}
\NormalTok{m }\SpecialCharTok{\%\textgreater{}\%} \FunctionTok{summarise}\NormalTok{(}\FunctionTok{across}\NormalTok{(}\FunctionTok{everything}\NormalTok{(), }\SpecialCharTok{\textasciitilde{}} \FunctionTok{sum}\NormalTok{(}\FunctionTok{is.na}\NormalTok{(.)))) }\SpecialCharTok{\%\textgreater{}\%}
  \FunctionTok{pivot\_longer}\NormalTok{(}\FunctionTok{everything}\NormalTok{(), }\AttributeTok{values\_to =} \StringTok{"n\_missing"}\NormalTok{) }\SpecialCharTok{\%\textgreater{}\%} \FunctionTok{filter}\NormalTok{(n\_missing }\SpecialCharTok{\textgreater{}} \DecValTok{0}\NormalTok{)}
\end{Highlighting}
\end{Shaded}

\begin{verbatim}
# A tibble: 5 x 2
  name       n_missing
  <chr>          <int>
1 ph                 1
2 nugent             2
3 amine_test         1
4 clue_cells         1
5 discharge          1
\end{verbatim}

\begin{Shaded}
\begin{Highlighting}[]
\FunctionTok{cat}\NormalTok{(}\StringTok{"duplicate sample\_ids:      "}\NormalTok{, }\FunctionTok{anyDuplicated}\NormalTok{(m}\SpecialCharTok{$}\NormalTok{sample\_id), }\StringTok{"}\SpecialCharTok{\textbackslash{}n}\StringTok{"}\NormalTok{)}
\end{Highlighting}
\end{Shaded}

\begin{verbatim}
duplicate sample_ids:       0 
\end{verbatim}

\begin{Shaded}
\begin{Highlighting}[]
\FunctionTok{cat}\NormalTok{(}\StringTok{"duplicate taxon rows:      "}\NormalTok{, }\FunctionTok{sum}\NormalTok{(}\FunctionTok{duplicated}\NormalTok{(d}\SpecialCharTok{$}\NormalTok{taxonomy}\SpecialCharTok{$}\NormalTok{taxon)), }\StringTok{"}\SpecialCharTok{\textbackslash{}n}\StringTok{"}\NormalTok{)}
\end{Highlighting}
\end{Shaded}

\begin{verbatim}
duplicate taxon rows:       0 
\end{verbatim}

\begin{Shaded}
\begin{Highlighting}[]
\FunctionTok{cat}\NormalTok{(}\StringTok{"duplicate gene rows:       "}\NormalTok{, }\FunctionTok{sum}\NormalTok{(}\FunctionTok{duplicated}\NormalTok{(d}\SpecialCharTok{$}\NormalTok{gene\_info}\SpecialCharTok{$}\NormalTok{ensembl\_gene\_id)), }\StringTok{"}\SpecialCharTok{\textbackslash{}n}\StringTok{"}\NormalTok{)}
\end{Highlighting}
\end{Shaded}

\begin{verbatim}
duplicate gene rows:        0 
\end{verbatim}

\begin{Shaded}
\begin{Highlighting}[]
\FunctionTok{cat}\NormalTok{(}\StringTok{"genes with no symbol:      "}\NormalTok{, }\FunctionTok{sum}\NormalTok{(}\FunctionTok{is.na}\NormalTok{(d}\SpecialCharTok{$}\NormalTok{gene\_info}\SpecialCharTok{$}\NormalTok{symbol)), }\StringTok{"}\SpecialCharTok{\textbackslash{}n}\StringTok{"}\NormalTok{)}
\end{Highlighting}
\end{Shaded}

\begin{verbatim}
genes with no symbol:       24 
\end{verbatim}

\begin{Shaded}
\begin{Highlighting}[]
\FunctionTok{cat}\NormalTok{(}\StringTok{"genes missing from counts: "}\NormalTok{,}
    \FunctionTok{length}\NormalTok{(}\FunctionTok{setdiff}\NormalTok{(d}\SpecialCharTok{$}\NormalTok{gene\_info}\SpecialCharTok{$}\NormalTok{ensembl\_gene\_id, }\FunctionTok{rownames}\NormalTok{(d}\SpecialCharTok{$}\NormalTok{genes))), }\StringTok{"}\SpecialCharTok{\textbackslash{}n}\StringTok{"}\NormalTok{)}
\end{Highlighting}
\end{Shaded}

\begin{verbatim}
genes missing from counts:  0 
\end{verbatim}

\subsection{7. Cytokines}\label{cytokines}

Measured concentrations, not counts, so the checks differ: gaps,
impossible negative values, and extreme outliers.

\begin{Shaded}
\begin{Highlighting}[]
\NormalTok{cy }\OtherTok{\textless{}{-}}\NormalTok{ d}\SpecialCharTok{$}\NormalTok{cytokines}
\FunctionTok{cat}\NormalTok{(}\StringTok{"samples:"}\NormalTok{, }\FunctionTok{nrow}\NormalTok{(cy), }\StringTok{" cytokines:"}\NormalTok{, }\FunctionTok{ncol}\NormalTok{(cy) }\SpecialCharTok{{-}} \DecValTok{1}\NormalTok{, }\StringTok{"}\SpecialCharTok{\textbackslash{}n}\StringTok{"}\NormalTok{)}
\end{Highlighting}
\end{Shaded}

\begin{verbatim}
samples: 352  cytokines: 18 
\end{verbatim}

\begin{Shaded}
\begin{Highlighting}[]
\FunctionTok{cat}\NormalTok{(}\StringTok{"missing values:  "}\NormalTok{, }\FunctionTok{sum}\NormalTok{(}\FunctionTok{is.na}\NormalTok{(cy)), }\StringTok{"}\SpecialCharTok{\textbackslash{}n}\StringTok{"}\NormalTok{)}
\end{Highlighting}
\end{Shaded}

\begin{verbatim}
missing values:   18 
\end{verbatim}

\begin{Shaded}
\begin{Highlighting}[]
\FunctionTok{cat}\NormalTok{(}\StringTok{"negative values: "}\NormalTok{, }\FunctionTok{sum}\NormalTok{(cy[, }\SpecialCharTok{{-}}\DecValTok{1}\NormalTok{] }\SpecialCharTok{\textless{}} \DecValTok{0}\NormalTok{, }\AttributeTok{na.rm =} \ConstantTok{TRUE}\NormalTok{), }\StringTok{"}\SpecialCharTok{\textbackslash{}n}\StringTok{"}\NormalTok{)}
\end{Highlighting}
\end{Shaded}

\begin{verbatim}
negative values:  0 
\end{verbatim}

\begin{Shaded}
\begin{Highlighting}[]
\CommentTok{\# Which cytokines have the gaps, and do they cluster in one arm?}
\FunctionTok{cat}\NormalTok{(}\StringTok{"missing values by cytokine:}\SpecialCharTok{\textbackslash{}n}\StringTok{"}\NormalTok{)}
\end{Highlighting}
\end{Shaded}

\begin{verbatim}
missing values by cytokine:
\end{verbatim}

\begin{Shaded}
\begin{Highlighting}[]
\NormalTok{count\_missing }\OtherTok{\textless{}{-}}\NormalTok{ cy }\SpecialCharTok{\%\textgreater{}\%}
  \FunctionTok{pivot\_longer}\NormalTok{(}\SpecialCharTok{{-}}\NormalTok{sample\_id) }\SpecialCharTok{\%\textgreater{}\%}
  \FunctionTok{group\_by}\NormalTok{(name) }\SpecialCharTok{\%\textgreater{}\%}
  \FunctionTok{summarise}\NormalTok{(}\AttributeTok{n\_missing =} \FunctionTok{sum}\NormalTok{(}\FunctionTok{is.na}\NormalTok{(value))) }\SpecialCharTok{\%\textgreater{}\%}
  \FunctionTok{filter}\NormalTok{(n\_missing }\SpecialCharTok{\textgreater{}} \DecValTok{0}\NormalTok{)}
\FunctionTok{print}\NormalTok{(count\_missing, }\AttributeTok{n =} \DecValTok{20}\NormalTok{)}
\end{Highlighting}
\end{Shaded}

\begin{verbatim}
# A tibble: 18 x 2
   name   n_missing
   <chr>      <int>
 1 ifng           1
 2 il_10          1
 3 il_13          1
 4 il_17          1
 5 il_1a          1
 6 il_1b          1
 7 il_21          1
 8 il_23          1
 9 il_4           1
10 il_6           1
11 il_8           1
12 ip_10          1
13 itac           1
14 mig            1
15 mip_1a         1
16 mip_1b         1
17 mip_3a         1
18 tnfa           1
\end{verbatim}

\begin{Shaded}
\begin{Highlighting}[]
\FunctionTok{cat}\NormalTok{(}\StringTok{"missing values by arm:}\SpecialCharTok{\textbackslash{}n}\StringTok{"}\NormalTok{)}
\end{Highlighting}
\end{Shaded}

\begin{verbatim}
missing values by arm:
\end{verbatim}

\begin{Shaded}
\begin{Highlighting}[]
\NormalTok{count\_missing\_arm }\OtherTok{\textless{}{-}}\NormalTok{ cy }\SpecialCharTok{\%\textgreater{}\%}
  \FunctionTok{left\_join}\NormalTok{(m }\SpecialCharTok{\%\textgreater{}\%} \FunctionTok{select}\NormalTok{(sample\_id, arm), }\AttributeTok{by =} \StringTok{"sample\_id"}\NormalTok{) }\SpecialCharTok{\%\textgreater{}\%}
  \FunctionTok{pivot\_longer}\NormalTok{(}\SpecialCharTok{{-}}\FunctionTok{c}\NormalTok{(sample\_id, arm)) }\SpecialCharTok{\%\textgreater{}\%}
  \FunctionTok{group\_by}\NormalTok{(arm) }\SpecialCharTok{\%\textgreater{}\%}
  \FunctionTok{summarise}\NormalTok{(}\AttributeTok{n\_missing =} \FunctionTok{sum}\NormalTok{(}\FunctionTok{is.na}\NormalTok{(value)))}
\FunctionTok{print}\NormalTok{(count\_missing\_arm, }\AttributeTok{n =} \DecValTok{5}\NormalTok{)}
\end{Highlighting}
\end{Shaded}

\begin{verbatim}
# A tibble: 2 x 2
  arm       n_missing
  <chr>         <int>
1 placebo           0
2 probiotic        18
\end{verbatim}

\begin{Shaded}
\begin{Highlighting}[]
\FunctionTok{summary}\NormalTok{(}\FunctionTok{as.matrix}\NormalTok{(cy[, }\SpecialCharTok{{-}}\DecValTok{1}\NormalTok{]))}
\end{Highlighting}
\end{Shaded}

\begin{verbatim}
      itac              mig                ip_10                ifng         
 Min.   : 0.7324   Min.   :    1.526   Min.   :   0.3052   Min.   :  0.3052  
 1st Qu.: 0.7324   1st Qu.:    5.598   1st Qu.:   3.3683   1st Qu.:  0.3052  
 Median : 1.7102   Median :   24.413   Median :   7.8235   Median :  1.0689  
 Mean   : 3.8398   Mean   :  725.412   Mean   :  92.3012   Mean   :  3.9787  
 3rd Qu.: 3.5416   3rd Qu.:  141.897   3rd Qu.:  37.4370   3rd Qu.:  3.0428  
 Max.   :50.4932   Max.   :12500.000   Max.   :2500.0000   Max.   :152.5871  
 NA's   :1         NA's   :1           NA's   :1           NA's   :1         
     il_17             il_13            il_21             il_8          
 Min.   : 0.3662   Min.   :0.1221   Min.   :0.1221   Min.   :   0.6013  
 1st Qu.: 0.3662   1st Qu.:0.1221   1st Qu.:0.1221   1st Qu.: 145.7826  
 Median : 0.3662   Median :0.1221   Median :0.1221   Median : 539.3313  
 Mean   : 2.7823   Mean   :0.3395   Mean   :0.3914   Mean   : 657.3246  
 3rd Qu.: 2.7368   3rd Qu.:0.4353   3rd Qu.:0.4323   3rd Qu.:1250.0000  
 Max.   :60.4507   Max.   :4.0863   Max.   :3.4495   Max.   :1250.0000  
 NA's   :1         NA's   :1        NA's   :1        NA's   :1          
     il_1a               il_1b               mip_1a             mip_1b        
 Min.   :   0.3052   Min.   :   0.2441   Min.   :  0.1526   Min.   :  0.4578  
 1st Qu.:  29.4669   1st Qu.:   0.8058   1st Qu.:  1.8950   1st Qu.:  0.4578  
 Median : 118.1873   Median :   5.3624   Median :  4.1969   Median :  1.5789  
 Mean   : 302.1808   Mean   :  68.0969   Mean   :  9.1473   Mean   :  6.4430  
 3rd Qu.: 363.9235   3rd Qu.:  35.2088   3rd Qu.:  8.6318   3rd Qu.:  5.6179  
 Max.   :2500.0000   Max.   :1635.3866   Max.   :312.6667   Max.   :176.4631  
 NA's   :1           NA's   :1           NA's   :1          NA's   :1         
     mip_3a              il_6               il_10              tnfa         
 Min.   :  0.3052   Min.   :  0.09155   Min.   : 0.7324   Min.   :  0.2136  
 1st Qu.:  1.3833   1st Qu.:  0.09155   1st Qu.: 0.7324   1st Qu.:  0.2136  
 Median :  3.2559   Median :  0.48833   Median : 0.7324   Median :  0.2136  
 Mean   : 16.3710   Mean   :  7.13916   Mean   : 2.1509   Mean   :  3.6276  
 3rd Qu.:  9.9805   3rd Qu.:  3.01529   3rd Qu.: 1.8213   3rd Qu.:  2.5276  
 Max.   :478.7761   Max.   :521.48895   Max.   :66.4982   Max.   :121.0065  
 NA's   :1          NA's   :1           NA's   :1         NA's   :1         
      il_4            il_23        
 Min.   :0.9155   Min.   :  3.967  
 1st Qu.:0.9155   1st Qu.:  3.967  
 Median :0.9155   Median :  9.771  
 Mean   :1.4971   Mean   : 24.947  
 3rd Qu.:1.9561   3rd Qu.: 26.262  
 Max.   :6.7704   Max.   :403.875  
 NA's   :1        NA's   :1        
\end{verbatim}

\begin{Shaded}
\begin{Highlighting}[]
\CommentTok{\# warning: false because geom\_histogram drops the 18 NA values; that is}
\CommentTok{\# intended here, and the gap itself is reported in the counts above.}
\NormalTok{cy }\SpecialCharTok{\%\textgreater{}\%}
  \FunctionTok{pivot\_longer}\NormalTok{(}\SpecialCharTok{{-}}\NormalTok{sample\_id) }\SpecialCharTok{\%\textgreater{}\%}
  \FunctionTok{ggplot}\NormalTok{(}\FunctionTok{aes}\NormalTok{(}\AttributeTok{x =}\NormalTok{ value, }\AttributeTok{fill =}\NormalTok{ name)) }\SpecialCharTok{+}
  \FunctionTok{geom\_histogram}\NormalTok{(}\AttributeTok{bins =} \DecValTok{40}\NormalTok{, }\AttributeTok{colour =} \StringTok{"white"}\NormalTok{, }\AttributeTok{linewidth =}\NormalTok{ .}\DecValTok{1}\NormalTok{, }\AttributeTok{alpha =}\NormalTok{ .}\DecValTok{9}\NormalTok{) }\SpecialCharTok{+}
  \FunctionTok{facet\_wrap}\NormalTok{(}\SpecialCharTok{\textasciitilde{}}\NormalTok{ name, }\AttributeTok{scales =} \StringTok{"free"}\NormalTok{) }\SpecialCharTok{+}
  \FunctionTok{scale\_fill\_hue}\NormalTok{(}\AttributeTok{l =} \DecValTok{45}\NormalTok{, }\AttributeTok{c =} \DecValTok{70}\NormalTok{) }\SpecialCharTok{+}
  \FunctionTok{labs}\NormalTok{(}\AttributeTok{title =} \StringTok{"Cytokine distributions (all 18)"}\NormalTok{, }\AttributeTok{x =} \StringTok{"concentration"}\NormalTok{, }\AttributeTok{y =} \StringTok{"samples"}\NormalTok{) }\SpecialCharTok{+}
  \FunctionTok{theme\_minimal}\NormalTok{(}\AttributeTok{base\_size =} \DecValTok{11}\NormalTok{) }\SpecialCharTok{+}
  \FunctionTok{theme}\NormalTok{(}\AttributeTok{legend.position =} \StringTok{"none"}\NormalTok{, }\AttributeTok{strip.text =} \FunctionTok{element\_text}\NormalTok{(}\AttributeTok{size =} \DecValTok{7}\NormalTok{))}
\end{Highlighting}
\end{Shaded}

\pandocbounded{\includegraphics[keepaspectratio]{01-qc_files/figure-pdf/cytokine-plot-1.pdf}}

\subsection{8. Verdict}\label{verdict}

\textbf{No blocking problems found.} The tables line up, and the count
tables are whole, non-negative and complete. The data is safe to take
into analysis.

Notes for the analysis stage:

\begin{itemize}
\tightlist
\item
  18 missing cytokine values
\item
  2 missing Nugent scores
\item
  microbiome depth outliers below a quarter of the median
\end{itemize}

\subsection{Choices made}\label{choices-made}

\begin{itemize}
\tightlist
\item
  Checked \textbf{raw counts} for integer/non-negative values; did not
  CPM-transform or log anything (transforming before DESeq2 is invalid).
\item
  Relied on \texttt{read\_counts()} to reorder count columns to
  metadata, but verified the order with \texttt{identical()};
  alternative would have been to reorder by hand.
\item
  Defined ``low-prevalence'' as present in fewer than 10\% of samples,
  as a number to report only --- no features were dropped here; the
  alternative of dropping them now was left to the DESeq2 step.
\item
  Checked library sizes by arm and time point visually; alternative
  would have been a formal test, which is overkill at QC stage.
\item
  Treated missing Nugent/post values as a data-gap note, not a failure;
  alternative was to exclude those people now, which would be premature
  before we know which analyses need Nugent.
\item
  Used ≥ 10\% prevalence and the default \texttt{summary()} ranges
  rather than setting custom thresholds, to keep QC reproducible without
  arbitrary cut-offs.
\end{itemize}




\end{document}
